\documentclass[11pt]{article}
\usepackage[a4paper,margin=1in]{geometry}
\usepackage{amsmath,amssymb}
\usepackage{hyperref}

\newcommand{\sgn}{\operatorname{sgn}}
\newcommand{\PV}{\mathrm{P.V.}\!\!}
\newcommand{\Hil}{\mathcal{H}}
\newcommand{\dd}{\mathrm{d}}
\newcommand{\vt}{v_t}

\title{The Hammett--Perkins closures:\\ where the Hilbert transform comes from}
\author{Working note}
\date{16 November 2026}

\begin{document}
\maketitle

\noindent\emph{Status.} This is a hand derivation written from first
principles. The closures, coefficients, response functions, and the
real-space integral \eqref{eq:real3} agree with a literature check against
Hammett--Perkins \cite{HP90} (eqs.~7--8), Hammett--Dorland--Perkins
\cite{HDP92}, and the review of Hunana et al.\ \cite{Hunana19}. The
intermediate algebra is my own and was not compared line by line; the
original papers are scans and some coefficients were confirmed through the
review and not read directly from them. See the convention notes in
Section~\ref{sec:caveat}.

\section{Setup and conventions}

Consider one species in one spatial dimension $z$ (the direction along the
magnetic field), with no collisions, linearized about a uniform Maxwellian
\[
  f_0(v) = \frac{n_0}{\sqrt{2\pi}\,\vt}\,
           \exp\!\left(-\frac{v^2}{2\vt^2}\right),
  \qquad \vt = \sqrt{T_0/m}.
\]
With acceleration $a = (e/m)E$, the perturbation $f_1$ obeys
\begin{equation}
  \partial_t f_1 + v\,\partial_z f_1 + a\,\partial_v f_0 = 0 .
  \label{eq:vlasov}
\end{equation}
Derivations below use units $n_0 = m = \vt = 1$; dimensional forms are
restored for the final results. Fourier modes are $\propto e^{ikz - i\omega t}$,
and
\begin{equation}
  \zeta = \frac{\omega}{\sqrt{2}\,|k|\,\vt}.
  \label{eq:zeta}
\end{equation}
The Hilbert transform is taken with the convention
\begin{equation}
  \Hil[g](z) = \frac{1}{\pi}\,\PV\int_{-\infty}^{\infty}
               \frac{g(z')}{z - z'}\,\dd z',
  \qquad
  \widehat{\Hil[g]}(k) = -i\,\sgn(k)\,\hat g(k).
  \label{eq:hilbert}
\end{equation}

\section{The kinetic response and the origin of $\sgn(k)$}

In Fourier space \eqref{eq:vlasov} gives, using $\partial_v f_0 = -v f_0$,
\[
  f_{1,k} = -\,\frac{i\,a\,v\,f_0(v)}{k v - \omega},
\]
so the density perturbation is
\[
  n_k = -\frac{i a}{k}\int \frac{v\,f_0}{v - \omega/k}\,\dd v
      = -\frac{i a}{k}\left[1 + \frac{\omega}{k}
        \int \frac{f_0}{v - \omega/k}\,\dd v\right].
\]
The integral is defined for $\operatorname{Im}\omega > 0$ (causality) and
continued analytically. It is the plasma dispersion function
\begin{equation}
  Z(\zeta) = \frac{1}{\sqrt{\pi}}\int_{-\infty}^{\infty}
             \frac{e^{-x^2}}{x - \zeta}\,\dd x ,
  \qquad \operatorname{Im}\zeta > 0,
  \label{eq:Z}
\end{equation}
which is the Hilbert transform, in velocity space, of a Gaussian. For real
$\zeta$ the Plemelj formula
\[
  \frac{1}{x - \zeta - i0} = \PV\frac{1}{x - \zeta} + i\pi\,\delta(x - \zeta)
\]
separates a principal-value part from the contribution of the resonant
particles at $v = \omega/k$.

The pole sits at $v = \omega/k$. For $\operatorname{Im}\omega > 0$ it lies
above the real $v$ axis when $k > 0$ and below it when $k < 0$. Substituting
$x \to -x$ in the second case shows that both signs give the same function of
$\zeta$ as defined in \eqref{eq:zeta}, with $|k|$ and not $k$:
\begin{equation}
  n_k = \frac{a}{ik}\,R(\zeta), \qquad R(\zeta) = 1 + \zeta Z(\zeta).
  \label{eq:R}
\end{equation}
This is the source of the non-analytic dependence on $k$. The expansions
needed below are
\begin{align}
  R(\zeta) &= 1 + i\sqrt{\pi}\,\zeta - 2\zeta^2 + O(\zeta^3),
  & \zeta &\to 0, \label{eq:Rsmall}\\
  R(\zeta) &= -\frac{1}{2\zeta^2} - \frac{3}{4\zeta^4}
              - \frac{15}{8\zeta^6} - \dots,
  & |\zeta| &\to \infty. \label{eq:Rlarge}
\end{align}
The term $i\sqrt{\pi}\,\zeta$ in \eqref{eq:Rsmall} is the pure resonant
contribution, $Z(0) = i\sqrt{\pi}$.

\section{Linearized moment equations}

Define $n = \int f_1\,\dd v$, $u = \int v f_1\,\dd v$,
$p = \int v^2 f_1\,\dd v$, $T = p - n$, the linearized heat flux
$q = \int v^3 f_1\,\dd v - 3u$, and the deviation of the fourth moment from
its Maxwellian value, $\delta r = r - 3p^2/n$, which linearizes to
$\delta r = \int v^4 f_1\,\dd v - 6p + 3n$. Moments of \eqref{eq:vlasov} give
\begin{align}
  \partial_t n + \partial_z u &= 0, \label{eq:n}\\
  \partial_t u + \partial_z (n + T) &= a, \label{eq:u}\\
  \partial_t T + 2\,\partial_z u + \partial_z q &= 0, \label{eq:T}\\
  \partial_t q + 3\,\partial_z T + \partial_z \delta r &= 0. \label{eq:q}
\end{align}
Each equation is exact; the hierarchy is unclosed because each involves the
next moment. In Fourier space write $\omega/k = \sqrt{2}\,\sgn(k)\,\zeta$.

\section{Three-moment closure}

Keep \eqref{eq:n}--\eqref{eq:T} and close $q$. A diffusive closure is
$q_k = -\chi\, i k\, T_k$. With collisions $\chi \sim \vt^2/\nu$; without
them the only available rate is the phase-mixing rate $|k|\vt$, so
$\chi = \chi_1 \sqrt{2}\,\vt/|k|$ with $\chi_1$ a pure number:
\begin{equation}
  q_k = -\,i\,\sgn(k)\,\sqrt{2}\,\chi_1\,T_k .
  \label{eq:closure3}
\end{equation}
Equation \eqref{eq:T} then gives $T_k = 2\zeta\,n_k/(\zeta + i\chi_1)$, and
\eqref{eq:u} gives $n_k = (a/ik)\,R_3$ with
\begin{equation}
  R_3(\zeta) = \frac{\chi_1 - i\zeta}
                    {\chi_1 - 3i\zeta - 2\chi_1\zeta^2 + 2i\zeta^3}.
  \label{eq:R3}
\end{equation}
For large $\zeta$, $R_3 = -1/(2\zeta^2) - 3/(4\zeta^4) + O(\zeta^{-5})$,
which agrees with \eqref{eq:Rlarge} for any $\chi_1$ because the retained
moment equations are exact. For small $\zeta$,
$R_3 = 1 + 2i\zeta/\chi_1 + O(\zeta^2)$; comparison with \eqref{eq:Rsmall}
fixes
\begin{equation}
  \chi_1 = \frac{2}{\sqrt{\pi}} .
\end{equation}
Restoring units,
\begin{equation}
  q_k = -\,n_0\,\chi_1\,\frac{\sqrt{2}\,\vt}{|k|}\,i k\,T_k
      = -\,i\,\sgn(k)\; n_0\sqrt{\frac{8}{\pi}}\;\vt\;T_k .
\end{equation}
By \eqref{eq:hilbert} the multiplier $-i\,\sgn(k)$ is the Hilbert transform,
so in real space
\begin{equation}
  q(z) = n_0\sqrt{\frac{8}{\pi}}\;\vt\;\Hil[T](z)
       = -\,n_0\,\chi_1\,\frac{\sqrt{2}\,\vt}{\pi}
         \int_0^\infty \frac{T(z+z') - T(z-z')}{z'}\,\dd z' .
  \label{eq:real3}
\end{equation}
The kernel carries no length scale, because a collisionless plasma has no
mean free path to supply one.

\section{Four-moment closure}

Now keep \eqref{eq:q} as well and close $\delta r$ in terms of the two
highest retained moments. Parity and dimensions allow a local term in $T$ and
a phase-mixing term in $q$:
\begin{equation}
  \delta r_k = \beta_1\,T_k - i\,\sgn(k)\,\sqrt{2}\,D_1\,q_k .
  \label{eq:closure4}
\end{equation}
The $T$ term is not dissipative; it replaces the coefficient $3$ in
\eqref{eq:q} by
\[
  b \equiv 3 + \beta_1 .
\]
Eliminating upward through the hierarchy:
\begin{align*}
  \text{from \eqref{eq:q}:}\quad
    q_k &= \frac{b\,T_k}{\sqrt{2}\,\sgn(k)\,(\zeta + iD_1)},\\[4pt]
  \text{from \eqref{eq:T}:}\quad
    T_k &= \frac{4\zeta(\zeta + iD_1)}{2\zeta(\zeta + iD_1) - b}\;n_k,\\[4pt]
  \text{from \eqref{eq:u}:}\quad
    n_k &= \frac{a}{ik}\,R_4(\zeta),
\end{align*}
with
\begin{equation}
  R_4(\zeta) = \frac{b - 2iD_1\zeta - 2\zeta^2}
                    {b - 6iD_1\zeta - (2b+6)\zeta^2
                     + 4iD_1\zeta^3 + 4\zeta^4}.
  \label{eq:R4}
\end{equation}
As in the three-moment case every factor of $\sgn(k)$ cancels except through
$\zeta$.

\paragraph{Matching at small $\zeta$.} With $\alpha = 2iD_1/b$,
\[
  R_4(\zeta) = 1 + 2\alpha\,\zeta
             + \left(2 + \frac{4}{b} + 6\alpha^2\right)\zeta^2 + O(\zeta^3).
\]
Comparing with \eqref{eq:Rsmall}:
\begin{align*}
  O(\zeta):&\quad 2\alpha = i\sqrt{\pi}
    \;\Longrightarrow\; D_1 = \frac{\sqrt{\pi}}{4}\,b,\\
  O(\zeta^2):&\quad 2 + \frac{4}{b} - \frac{3\pi}{2} = -2
    \;\Longrightarrow\; b = \frac{8}{3\pi - 8}.
\end{align*}
Therefore
\begin{equation}
  D_1 = \frac{2\sqrt{\pi}}{3\pi - 8} \approx 2.49,
  \qquad
  \beta_1 = \frac{32 - 9\pi}{3\pi - 8} \approx 2.62 .
  \label{eq:coeffs4}
\end{equation}

\paragraph{Behaviour at large $\zeta$.} From \eqref{eq:R4},
\[
  \frac{1}{R_4} = -2\zeta^2 + 3 + \frac{b}{\zeta^2} + O(\zeta^{-3}),
  \qquad\text{whereas}\qquad
  \frac{1}{R} = -2\zeta^2 + 3 + \frac{3}{\zeta^2} + \dots
\]
So $R_4$ reproduces the first two terms of \eqref{eq:Rlarge} for any $D_1$
and $\beta_1$, but the $\zeta^{-6}$ term is correct only if $\beta_1 = 0$.
The closure spends its two parameters on the low-frequency side and gives up
one order of high-frequency accuracy to do so.

\paragraph{Real space.} Restoring units, \eqref{eq:closure4} reads
\begin{equation}
  \delta r_k = \beta_1\,n_0\vt^2\,T_k
             - D_1\,\frac{\sqrt{2}\,\vt}{|k|}\,i k\,q_k ,
  \qquad
  \delta r(z) = \beta_1\,n_0\vt^2\,T(z)
              + \sqrt{2}\,D_1\,\vt\,\Hil[q](z),
\end{equation}
to be used in the dimensional form of \eqref{eq:q},
$\partial_t q + 3 n_0\vt^2\,\partial_z T + \partial_z\delta r = 0$.
The Hilbert transform now acts on the heat flux, which has its own evolution
equation, and not on the temperature.

\section{Summary}

\begin{center}
\renewcommand{\arraystretch}{1.5}
\begin{tabular}{lll}
  & Three-moment & Four-moment \\ \hline
  Closed quantity & $q$ & $\delta r$ \\
  Closure & $q = n_0\sqrt{2}\,\chi_1\vt\,\Hil[T]$
          & $\delta r = \beta_1 n_0\vt^2 T + \sqrt{2}\,D_1\vt\,\Hil[q]$ \\
  Coefficients & $\chi_1 = 2/\sqrt{\pi}$
          & $D_1 = \dfrac{2\sqrt{\pi}}{3\pi-8}$,\;
            $\beta_1 = \dfrac{32-9\pi}{3\pi-8}$ \\
  Matched as $\zeta \to 0$ & through $O(\zeta)$ & through $O(\zeta^2)$ \\
  Matched as $\zeta \to \infty$ & through $O(\zeta^{-4})$
          & through $O(\zeta^{-4})$ \\
\end{tabular}
\end{center}

In both cases the response function is a rational (Pad\'e-type) approximant
of $R(\zeta) = 1 + \zeta Z(\zeta)$, and the Hilbert transform in $z$ is the
image of the Landau resonance: the resonant denominator makes the response a
function of $\omega/|k|$, the closure coefficient scales as $1/|k|$, and
$(1/|k|)\,ik = i\,\sgn(k)$.

\section{Caveats}
\label{sec:caveat}

\begin{itemize}
  \item \textbf{Thermal-speed convention.} This note and
    Hammett--Perkins use $\vt = \sqrt{T_0/m}$, for which the coefficients are
    those in \eqref{eq:coeffs4}. Hunana et al.\ \cite{Hunana19} use
    $v_{\mathrm{th}} = \sqrt{2T_0/m}$, in which the same closure reads
    \[
      \delta r_k = -\,i\,\frac{2\sqrt{\pi}}{3\pi - 8}\,v_{\mathrm{th}}\,
                   \sgn(k)\,q_k
                 + \frac{32 - 9\pi}{2(3\pi - 8)}\,v_{\mathrm{th}}^2\,n_0\,T_k ,
    \]
    so the local coefficient appears a factor of two smaller, and the
    three-moment closure reads
    $q_k = -i\,(2/\sqrt{\pi})\,n_0 v_{\mathrm{th}}\,\sgn(k)\,T_k$.
  \item \textbf{Sign of the Hilbert transform.} With \eqref{eq:hilbert} and
    modes $e^{+ikz}$, $i\,\sgn(k) \leftrightarrow -\Hil$. Several
    Landau-fluid papers define the opposite sign so that
    $i\,\sgn(k) \leftrightarrow \Hil$; the integral form \eqref{eq:real3}
    is convention-free. The 1990 and 1992 papers write the integral but do
    not use the words ``Hilbert transform''; the identification is made
    explicitly in \cite{Hunana19}.
  \item \textbf{Known inaccuracies.} According to \cite{Hunana19}, the
    three-moment closure gives $q \sim \zeta^{-2}$ at large $\zeta$ where
    the kinetic result is $\zeta^{-3}$, and so tends to overestimate Landau
    damping; the closures are also local in time.
  \item \textbf{Validity.} The matching is linear, about a Maxwellian, and
    exact only in the limits $\zeta \to 0$ and $\zeta \to \infty$; in between
    the closures are approximants.
  \item \textbf{Domain.} $\Hil$ is defined on the whole line or a periodic
    domain. Bounded or sheath-terminated field lines require a separate
    choice of extension.
\end{itemize}

\begin{thebibliography}{9}
\bibitem{HP90}
G.~W. Hammett and F.~W. Perkins,
\emph{Fluid moment models for Landau damping with application to the
ion-temperature-gradient instability},
Phys.\ Rev.\ Lett.\ \textbf{64}, 3019 (1990).
\bibitem{HDP92}
G.~W. Hammett, W.~Dorland, and F.~W. Perkins,
\emph{Fluid models of phase mixing, Landau damping, and nonlinear
gyrokinetic dynamics},
Phys.\ Fluids B \textbf{4}, 2052 (1992).
\bibitem{Hunana19}
P.~Hunana et al.,
\emph{An introductory guide to fluid models with anisotropic temperatures.
Part 2. Kinetic theory, Pad\'e approximants and Landau fluid closures},
\href{https://arxiv.org/abs/1901.09360}{arXiv:1901.09360} (2019).
\end{thebibliography}

\end{document}
